\section{Related Work}

\label{sec:related}

\subsection{Retrieval-Augmented Generation}

RAG was introduced by Lewis et al.~\cite{lewis2020retrieval} as a framework for combining retrieval with generation. REALM~\cite{guu2020realm} pre-trains a retriever jointly with a language model. Atlas~\cite{izacard2023atlas} achieves strong few-shot performance with retrieval. Self-RAG~\cite{asai2023self} adds self-reflection tokens for adaptive retrieval. REPLUG~\cite{shi2023replug} treats the retriever as a plug-in. Hanzo Search extends RAG with hybrid retrieval, sentence-level attribution, and generative UI.

\subsection{Generative Search Engines}

Perplexity AI~\cite{perplexity2023} pioneered the conversational search format. Bing Copilot~\cite{nair2023bing} integrates GPT-4 with Bing search. Google SGE~\cite{google2023sge} adds AI summaries to Google Search. You.com~\cite{you2023} combines search with AI chat. Hanzo Search differentiates through generative UI components, learned hybrid retrieval, and sentence-level attribution.

\subsection{Hybrid Retrieval}

RRF~\cite{cormack2009reciprocal} provides a simple rank fusion method. ColBERTv2~\cite{santhanam2022colbertv2} uses late interaction for efficient dense retrieval. SPLADE~\cite{formal2022splade} learns sparse representations. Hanzo Search uses learned fusion weights conditioned on query features, achieving better adaptation than fixed-weight or rank-based methods.

\subsection{Source Attribution}

ALCE~\cite{gao2023alce} provides benchmarks for attribution in LLM generation. RARR~\cite{gao2023rarr} retrofits attribution to existing responses. WebGPT~\cite{nakano2021webgpt} trains models to cite sources. Hanzo Search implements inline citation injection with post-hoc NLI verification.

